\documentclass[11pt,a4paper]{article}
\usepackage[top=2.3cm,bottom=2.3cm,left=2.2cm,right=2.2cm]{geometry}

\usepackage{amsmath,amssymb}
\usepackage{fontspec}
\setmainfont{Liberation Serif}
\setsansfont{Liberation Sans}
\setmonofont{DejaVu Sans Mono}[Scale=0.84]

\usepackage{xcolor}
\definecolor{brandblue}{RGB}{20, 55, 105}
\definecolor{accentblue}{RGB}{35, 95, 165}
\definecolor{codebg}{RGB}{248, 249, 251}
\definecolor{codeframe}{RGB}{215, 222, 230}
\definecolor{javakeyword}{RGB}{0, 45, 170}
\definecolor{javastring}{RGB}{5, 120, 20}
\definecolor{javacomment}{RGB}{120, 120, 120}
\definecolor{javanumber}{RGB}{20, 75, 220}

\definecolor{noteborder}{RGB}{30, 110, 65}
\definecolor{notebg}{RGB}{242, 249, 245}
\definecolor{warnborder}{RGB}{185, 75, 10}
\definecolor{warnbg}{RGB}{254, 248, 240}
\definecolor{tipborder}{RGB}{30, 95, 160}
\definecolor{tipbg}{RGB}{242, 247, 254}

\usepackage{fancyhdr}
\usepackage{lastpage}
\pagestyle{fancy}
\fancyhf{}
\lhead{\parbox[t]{0.58\textwidth}{\raggedright\small\sffamily\color{brandblue}\textbf{T13 --- Classi Astratte ed Interfacce}}}
\rhead{\small\sffamily\color{gray}Programmazione II -- UniVR (2025/2026)}
\cfoot{\small\sffamily Pagina \thepage\ di \pageref{LastPage}}
\renewcommand{\headrulewidth}{0.5pt}
\renewcommand{\headrule}{\hbox to\headwidth{\color{brandblue!70!black}\leaders\hrule height \headrulewidth\hfill}}

\usepackage{titlesec}
\titleformat{\section}{\Large\bfseries\sffamily\color{brandblue}}{\thesection}{0.8em}{}[\titlerule]
\titleformat{\subsection}{\large\bfseries\sffamily\color{accentblue}}{\thesubsection}{0.8em}{}
\titleformat{\subsubsection}{\normalsize\bfseries\sffamily\color{brandblue!85!black}}{\thesubsubsection}{0.8em}{}

\usepackage{needspace}
\usepackage{fancyvrb}
\DefineVerbatimEnvironment{asciibox}{Verbatim}{
  fontsize=\fontsize{7.5pt}{9.2pt}\selectfont,
  frame=single,
  rulecolor=\color{codeframe},
  fillcolor=\color{codebg},
  framesep=2.5mm
}
\usepackage{listings}
\lstset{
  backgroundcolor=\color{codebg},
  frame=single,
  rulecolor=\color{codeframe},
  basicstyle=\ttfamily\footnotesize,
  keywordstyle=\color{javakeyword}\bfseries,
  stringstyle=\color{javastring},
  commentstyle=\color{javacomment}\itshape,
  numberstyle=\tiny\color{gray},
  numbers=left,
  numbersep=7pt,
  stepnumber=1,
  breaklines=true,
  breakatwhitespace=true,
  showstringspaces=false,
  tabsize=4,
  xleftmargin=12pt,
  framexleftmargin=12pt
}

\newcommand{\passthrough}[1]{#1}
\providecommand{\tightlist}{\setlength{\itemsep}{0pt}\setlength{\parskip}{0pt}}

\usepackage{tcolorbox}
\tcbuselibrary{skins,breakable}
\newtcolorbox{studynote}[1][Nota]{
  enhanced,
  breakable,
  colback=notebg,
  colframe=noteborder,
  fonttitle=\bfseries\sffamily\small,
  coltitle=white,
  title={#1},
  arc=2mm,
  boxrule=0.8pt,
  left=9pt, right=9pt, top=7pt, bottom=7pt
}

\newtcolorbox{studywarning}[1][Attenzione]{
  enhanced,
  breakable,
  colback=warnbg,
  colframe=warnborder,
  fonttitle=\bfseries\sffamily\small,
  coltitle=white,
  title={#1},
  arc=2mm,
  boxrule=0.8pt,
  left=9pt, right=9pt, top=7pt, bottom=7pt
}

\newtcolorbox{studytip}[1][Suggerimento]{
  enhanced,
  breakable,
  colback=tipbg,
  colframe=tipborder,
  fonttitle=\bfseries\sffamily\small,
  coltitle=white,
  title={#1},
  arc=2mm,
  boxrule=0.8pt,
  left=9pt, right=9pt, top=7pt, bottom=7pt
}

\usepackage{booktabs}
\usepackage{longtable}
\usepackage{array}
\usepackage{calc}

\usepackage{hyperref}
\hypersetup{
  colorlinks=true,
  linkcolor=accentblue,
  urlcolor=accentblue,
  citecolor=brandblue,
  pdfborder={0 0 0}
}

\title{\Huge\bfseries\sffamily\color{brandblue} T13 --- Classi Astratte ed Interfacce \\[0.3em] \large\sffamily Contratti Comportamentali, Metodi Astratti, Default e Static Methods, Builder Pattern e Lambda}
\author{\textbf{Docente:} Prof. Michele Pasqua \quad---\quad \textbf{Appunti Revisione:} Wally}
\date{A.A. 2025/2026 -- Universit\`a degli Studi di Verona}

\begin{document}
\maketitle
\thispagestyle{fancy}
\vspace{-1.2em}
\noindent\textcolor{brandblue}{\rule{\textwidth}{0.8pt}}
\vspace{1.2em}

In questo blocco analizziamo a fondo i meccanismi di astrazione pura di
Java: le \textbf{Classi Astratte}, le \textbf{Interfacce}, il design
pattern \textbf{Template Method} e le interfacce cardine della libreria
standard (\passthrough{\lstinline!Comparable!},
\passthrough{\lstinline!Iterable!}, \passthrough{\lstinline!Iterator!})
(\href{file:///home/wally/Documenti/GitHub/Programmazione-II/25-26/Teoria-e-Esercitazioni/Slides-20260902/T13\%20-\%20Abstract\%20Classes\%20and\%20Interfaces.pdf}{T13
- Abstract Classes and Interfaces.pdf}). Nel codice di
\href{file:///home/wally/Documenti/GitHub/Programmazione-II/25-26/Teoria-e-Esercitazioni/Codice-20260902/Lezione13/MavenDate}{\passthrough{\lstinline!Lezione13/MavenDate!}},
il docente compie un'evoluzione straordinaria introducendo: 1. La classe
statica annidata
\textbf{\href{file:///home/wally/Documenti/GitHub/Programmazione-II/25-26/Teoria-e-Esercitazioni/Codice-20260902/Lezione13/MavenDate/src/main/java/it/oop/core/Date.java}{\passthrough{\lstinline!Date.Builder!}}}
per la costruzione controllata degli oggetti. 2. L'interfaccia
funzionale
\textbf{\href{file:///home/wally/Documenti/GitHub/Programmazione-II/25-26/Teoria-e-Esercitazioni/Codice-20260902/Lezione13/MavenDate/src/main/java/it/oop/core/FormattedDateConverter.java}{\passthrough{\lstinline!FormattedDateConverter!}}}.
3. L'uso congiunto in
\href{file:///home/wally/Documenti/GitHub/Programmazione-II/25-26/Teoria-e-Esercitazioni/Codice-20260902/Lezione13/MavenDate/src/main/java/it/oop/ui/MainDate.java}{\passthrough{\lstinline!MainDate.java!}}
di \textbf{Classi Anonime}
(\passthrough{\lstinline!new Time() \{ ... \}!}) e \textbf{Espressioni
Lambda} (\passthrough{\lstinline!d -> new AmericanDate(...)!}).

\begin{center}\rule{0.5\linewidth}{0.5pt}\end{center}

\subsubsection{1. Teoria: Concetti Teorici dalle Slide (Slide
T13)}\label{teoria-concetti-teorici-dalle-slide-slide-t13}

\paragraph{A. Classi Astratte e Template Method Pattern (Slide
1--10)}\label{a.-classi-astratte-e-template-method-pattern-slide-110}

\begin{itemize}
\tightlist
\item
  \textbf{Metodo Astratto}:

  \begin{itemize}
  \item
    Un metodo dichiarato con la parola chiave
    \passthrough{\lstinline!abstract!} e \textbf{privo di corpo}
    (termina con il punto e virgola \passthrough{\lstinline!;!} anziché
    con le parentesi graffe \passthrough{\lstinline!\{ ... \}!}):

\begin{lstlisting}[language=Java]
public abstract String prettyPrint();
\end{lstlisting}
  \item
    Definisce una ``firma obbligatoria'': impone alle sottoclassi
    concrete il compito di implementarne la logica.
  \end{itemize}
\item
  \textbf{Classe Astratta (\passthrough{\lstinline!abstract class!})}:

  \begin{itemize}
  \tightlist
  \item
    Se una classe contiene anche un solo metodo astratto, \textbf{deve
    obbligatoriamente essere dichiarata
    \passthrough{\lstinline!abstract!}}.
  \item
    Attenzione: \textbf{Regola aurea d'esame}: Una classe astratta
    \textbf{non può MAI essere istanziata direttamente}
    (\passthrough{\lstinline!new FormattedDate(...)!} genera un
    \textbf{Compile-Time Error}!).
  \item
    Può contenere costruttori (invocabili dalle sottoclassi tramite
    \passthrough{\lstinline!super(...)!}), campi di istanza
    (\passthrough{\lstinline!protected!}/\passthrough{\lstinline!private!}),
    metodi concreti e metodi astratti.
  \item
    Una sottoclasse che estende una classe astratta deve:

    \begin{itemize}
    \tightlist
    \item
      O implementare \textbf{tutti} i metodi astratti ereditati.
    \item
      Oppure essere dichiarata essa stessa
      \passthrough{\lstinline!abstract!}.
    \end{itemize}
  \end{itemize}
\item
  \textbf{Il Template Method Pattern}:

  \begin{itemize}
  \tightlist
  \item
    Pattern comportamentale in cui la superclasse definisce lo scheletro
    immutabile di un algoritmo all'interno di un metodo concreto marcato
    \textbf{\passthrough{\lstinline!final!}} (il template), mentre
    delega i singoli passi variabili o dipendenti dal contesto a metodi
    \passthrough{\lstinline!abstract!} implementati dalle sottoclassi.
  \item
    Nel nostro progetto: \passthrough{\lstinline!printFormat()!} e
    \passthrough{\lstinline!getMonthAsString()!} sono
    \passthrough{\lstinline!final!} in
    \passthrough{\lstinline!FormattedDate!}, mentre
    \passthrough{\lstinline!prettyPrint()!} è
    \passthrough{\lstinline!abstract!}.
  \end{itemize}
\end{itemize}

\begin{center}\rule{0.5\linewidth}{0.5pt}\end{center}

\paragraph{\texorpdfstring{B. Le Interfacce in Java (\texttt{interface})
(Slide
11--17)}{B. Le Interfacce in Java (interface) (Slide 11--17)}}\label{b.-le-interfacce-in-java-interface-slide-1117}

\begin{itemize}
\tightlist
\item
  \textbf{Definizione e Filosofia}:

  \begin{itemize}
  \tightlist
  \item
    Un'interfaccia è un \textbf{contratto puro di comportamento}.
    Rappresenta ciò che una classe \emph{sa fare} (\emph{can-do}), non
    la sua identità ontologica.
  \item
    \textbf{Membri di un'interfaccia standard}:

    \begin{itemize}
    \tightlist
    \item
      \textbf{Campi}: sono implicitamente ed esclusivamente
      \textbf{\passthrough{\lstinline!public static final!}} (costanti
      di classe). Non possono esistere variabili d'istanza o stato
      mutabile.
    \item
      \textbf{Metodi}: sono implicitamente
      \textbf{\passthrough{\lstinline!public abstract!}} (non serve
      specificarlo).
    \end{itemize}
  \end{itemize}
\item
  \textbf{Ereditarietà Multipla di Tipo
  (\passthrough{\lstinline!implements!})}:

  \begin{itemize}
  \item
    Una classe può estendere al massimo una sola superclasse, ma può
    implementare \textbf{un numero arbitrario di interfacce separate da
    virgola}:

\begin{lstlisting}[language=Java]
public class TimeStamp extends Date implements Time, Serializable, Cloneable
\end{lstlisting}
  \item
    Questo realizza in Java l'ereditarietà multipla dei tipi senza
    incorrere nei problemi di ambiguità di memoria del C++.
  \end{itemize}
\item
  \textbf{Ereditarietà tra Interfacce}:

  \begin{itemize}
  \tightlist
  \item
    Un'interfaccia può estendere altre interfacce tramite
    \passthrough{\lstinline!extends!} (anche multiple
    contemporaneamente:
    \passthrough{\lstinline!interface C extends A, B!}).
  \end{itemize}
\end{itemize}

\begin{center}\rule{0.5\linewidth}{0.5pt}\end{center}

\paragraph{C. Confronto Sistematico: Classe Astratta vs
Interfaccia}\label{c.-confronto-sistematico-classe-astratta-vs-interfaccia}

\begin{longtable}[]{@{}
  >{\raggedright\arraybackslash}p{(\linewidth - 4\tabcolsep) * \real{0.3333}}
  >{\raggedright\arraybackslash}p{(\linewidth - 4\tabcolsep) * \real{0.3333}}
  >{\raggedright\arraybackslash}p{(\linewidth - 4\tabcolsep) * \real{0.3333}}@{}}
\toprule\noalign{}
\begin{minipage}[b]{\linewidth}\raggedright
Proprietà
\end{minipage} & \begin{minipage}[b]{\linewidth}\raggedright
Classe Astratta (\passthrough{\lstinline!abstract class!})
\end{minipage} & \begin{minipage}[b]{\linewidth}\raggedright
Interfaccia (\passthrough{\lstinline!interface!})
\end{minipage} \\
\midrule\noalign{}
\endhead
\bottomrule\noalign{}
\endlastfoot
\textbf{Istanziabilità} & {[}NO{]} No (\passthrough{\lstinline!new!}
vietato) & {[}NO{]} No (\passthrough{\lstinline!new!} vietato) \\
\textbf{Stato d'istanza (campi)} & {[}OK{]} Sì (qualsiasi visibilità) &
{[}NO{]} No (solo costanti
\passthrough{\lstinline!public static final!}) \\
\textbf{Costruttori} & {[}OK{]} Sì (invocabili con
\passthrough{\lstinline!super(...)!}) & {[}NO{]} No \\
\textbf{Ereditarietà} & Singola (\passthrough{\lstinline!extends!} una
sola classe) & Multipla (\passthrough{\lstinline!implements!} \(N\)
interfacce) \\
\textbf{Relazione concettuale} & Identità ontologica forte (\emph{IS-A})
& Capacità / Contratto comportamentale (\emph{CAN-DO}) \\
\end{longtable}

\begin{center}\rule{0.5\linewidth}{0.5pt}\end{center}

\paragraph{D. Interfacce Standard Fondamentali della JDK (Slide
18--22)}\label{d.-interfacce-standard-fondamentali-della-jdk-slide-1822}

\begin{enumerate}
\def\labelenumi{\arabic{enumi}.}
\tightlist
\item
  \textbf{\passthrough{\lstinline!java.lang.Comparable<T>!}}:

  \begin{itemize}
  \tightlist
  \item
    Definisce il metodo
    \passthrough{\lstinline!int compareTo(T other)!}.
  \item
    Permette agli algoritmi di ordinamento
    (\passthrough{\lstinline!Arrays.sort()!},
    \passthrough{\lstinline!Collections.sort()!}) e agli insiemi
    ordinati (\passthrough{\lstinline!TreeSet!},
    \passthrough{\lstinline!TreeMap!}) di ordinare gli oggetti secondo
    il loro ``ordine naturale''.
  \end{itemize}
\item
  \textbf{\passthrough{\lstinline!java.lang.Iterable<T>!} e
  \passthrough{\lstinline!java.util.Iterator<T>!}}:

  \begin{itemize}
  \tightlist
  \item
    \textbf{\passthrough{\lstinline!Iterable<T>!}}: espone il metodo
    \passthrough{\lstinline!Iterator<T> iterator()!}.
  \item
    \textbf{Regola cruciale}: Qualsiasi classe che implementi
    \passthrough{\lstinline!Iterable!} può essere utilizzata come
    sorgente nel \textbf{ciclo For-Each
    (\passthrough{\lstinline!for (T elem : collection)!})}!
  \item
    \textbf{\passthrough{\lstinline!Iterator<T>!}}: l'oggetto cursore
    che scorre la sequenza:

    \begin{itemize}
    \tightlist
    \item
      \passthrough{\lstinline!boolean hasNext()!}: verifica se vi sono
      ulteriori elementi.
    \item
      \passthrough{\lstinline!T next()!}: restituisce il prossimo
      elemento avanzando il cursore.
    \item
      \passthrough{\lstinline!default void remove()!}: rimuove l'ultimo
      elemento restituito.
    \end{itemize}
  \end{itemize}
\end{enumerate}

\begin{center}\rule{0.5\linewidth}{0.5pt}\end{center}

\subsubsection{\texorpdfstring{2. Codice: Evoluzione del Codice:
\href{file:///home/wally/Documenti/GitHub/Programmazione-II/25-26/Teoria-e-Esercitazioni/Codice-20260902/Lezione13/MavenDate}{\texttt{Lezione13/MavenDate}}}{2. Codice: Evoluzione del Codice: Lezione13/MavenDate}}\label{codice-evoluzione-del-codice-lezione13mavendate}

\paragraph{\texorpdfstring{1. Il Design Pattern Builder con Classe
Statica Annidata
(\href{file:///home/wally/Documenti/GitHub/Programmazione-II/25-26/Teoria-e-Esercitazioni/Codice-20260902/Lezione13/MavenDate/src/main/java/it/oop/core/Date.java}{\texttt{Date.java}})}{1. Il Design Pattern Builder con Classe Statica Annidata (Date.java)}}\label{il-design-pattern-builder-con-classe-statica-annidata-date.java}

All'interno di \passthrough{\lstinline!Date.java!}, il docente definisce
una \textbf{Static Nested Class} per costruire date in modo protetto:

\begin{lstlisting}[language=Java]
public class Date implements Comparable {
    ...
    public static class Builder {
        private final int year;

        public Builder(int year) {
            if (year > 0)
                this.year = year;
            else
                this.year = 1970; // Anno di default sicuro
        }

        public Date build(int day, int month) {
            // Se i parametri sono illegali, ricade sulla data sicura 1/1/year!
            if (month < 1 || month > 12 || day < 1 || day > daysPerMonth(month))
                return new Date(1, 1, year);
            return new Date(day, month, year);
        }
    }
}
\end{lstlisting}

\emph{Vantaggio del Builder}: Incapsula e pre-configura l'anno; se
l'utente fornisce parametri non validi
(\passthrough{\lstinline!month = -1!}), il costruttore non fallisce a
video ma restituisce una data valida di fallback
(\passthrough{\lstinline!01/01/year!}).

\begin{center}\rule{0.5\linewidth}{0.5pt}\end{center}

\paragraph{\texorpdfstring{2. L'Interfaccia Funzionale:
\href{file:///home/wally/Documenti/GitHub/Programmazione-II/25-26/Teoria-e-Esercitazioni/Codice-20260902/Lezione13/MavenDate/src/main/java/it/oop/core/FormattedDateConverter.java}{\texttt{FormattedDateConverter.java}}}{2. L'Interfaccia Funzionale: FormattedDateConverter.java}}\label{linterfaccia-funzionale-formatteddateconverter.java}

\begin{lstlisting}[language=Java]
package it.oop.core;

public interface FormattedDateConverter {
    FormattedDate convert(FormattedDate date);
}
\end{lstlisting}

Questa interfaccia dichiara \textbf{un solo metodo astratto} (SAM:
\emph{Single Abstract Method}). In Java è a tutti gli effetti
un'\textbf{Interfaccia Funzionale}, target ideale per le espressioni
Lambda!

\begin{center}\rule{0.5\linewidth}{0.5pt}\end{center}

\paragraph{\texorpdfstring{3. Sintesi Pratica in
\href{file:///home/wally/Documenti/GitHub/Programmazione-II/25-26/Teoria-e-Esercitazioni/Codice-20260902/Lezione13/MavenDate/src/main/java/it/oop/ui/MainDate.java}{\texttt{MainDate.java}}}{3. Sintesi Pratica in MainDate.java}}\label{sintesi-pratica-in-maindate.java}

Questo file è una miniera di concetti d'esame avanzati:

\begin{lstlisting}[language=Java]
package it.oop.ui;

import it.oop.core.*;
import it.oop.core.Date;

public class MainDate {
    public static void main(String[] args) {
        // 1. Uso della Static Nested Class Builder:
        Date.Builder dateBuilder = new Date.Builder(2025);
        Date d1 = dateBuilder.build(10, 9);  // Giorno 10, Mese 9 -> Valida!
        Date d2 = dateBuilder.build(10, -1); // Mese -1 illegale -> Fallback su 1/1/2025!
        System.out.println(d1.toString());   // y2025m9d10
        System.out.println(d2.toString());   // y2025m1d1

        // 2. Classe Anonima (Anonymous Inner Class):
        // Implementa al volo l'interfaccia Time senza creare un file .java separato!
        Time init = new Time() {
            @Override
            public int getHours() { return 0; }
            @Override
            public int getMinutes() { return 0; }
            @Override
            public int getSeconds() { return 0; }
            @Override
            public String toString() {
                return String.format("%02d:%02d:%02d", getHours(), getMinutes(), getSeconds());
            }
        };
        System.out.println(init.toString()); // Stampa 00:00:00

        // 3. Espressione Lambda (Sintassi compatta per Interfaccia Funzionale):
        FormattedDateConverter toAmerican =
                d -> new AmericanDate(d.getDay(), d.getMonth(), d.getYear());

        // Test di conversione polimorfica:
        System.out.println(toAmerican.convert(new ItalianDate(11, 11, 2025)) instanceof AmericanDate);
        // Stampa TRUE!
    }
}
\end{lstlisting}

\begin{center}\rule{0.5\linewidth}{0.5pt}\end{center}

\subsubsection{3. Test: Output di Esecuzione Verificato con
Maven}\label{test-output-di-esecuzione-verificato-con-maven}

Eseguendo
\passthrough{\lstinline!mvn exec:java -Dexec.mainClass="it.oop.ui.MainDate"!}:

\begin{lstlisting}
y2025m9d10
y2025m1d1
00:00:00
true
\end{lstlisting}

\begin{center}\rule{0.5\linewidth}{0.5pt}\end{center}

\begin{studynote}[Nota di Studio] Con la \textbf{Lezione 13} abbiamo chiuso il cerchio su
Classi Astratte, Interfacce, Template Method e abbiamo anticipato le
classi annidate e le lambda.

Il prossimo blocco è la \textbf{Lezione 14 / Slide T14: \emph{Nested and
Anonymous Classes}}: - Tassonomia completa delle classi interne:
\textbf{Static Nested Classes}, \textbf{Inner Member Classes} (non
statiche), \textbf{Local Classes} e \textbf{Anonymous Classes}. -
Accesso allo stato della classe contenitore (\emph{Enclosing instance})
e la sintassi speciale \passthrough{\lstinline!EnclosingClass.this!}. -
Cattura delle variabili locali nelle classi locali/anonime e il vincolo
di essere \textbf{\passthrough{\lstinline!effectively final!}}. -
Evoluzione del codice in \passthrough{\lstinline!Lezione14!}:
consolidamento delle classi interne e preparativi per i Generics.
\end{studynote}

\textbf{Dammi conferma per aprire la Lezione 14 / T14 e continuare la
revisione!}


\end{document}
